\documentclass{beamer}
%[aspectratio=169]   \usepackage[czech]{babel}
\usepackage{apo-lecture-cz}
\usepackage{pdfpages}
\usepackage{pdfcomment}
\usepackage{listings}
\usepackage{array,multirow}

\subtitle{Lekce 02. Reprezentace čísel}
\author{Petr Štěpán \phantom{xxxxxxx} Pavel Píša \\ \small\texttt{stepan@fel.cvut.cz}\phantom{xxxx}\small\texttt{pisa@fel.cvut.cz}}
\begin{document}

\maketitle

\section{ALU - dokončení}

\begin{frame}
\frametitle{Násobení celých čísel}

Podle uvedeného algoritmu můžeme vytvořit následující násobičku s posuvným registrem:\\
(A,B 32 bitů, výsledek 64 bitů)
\begin{center}
\includegraphics[width=0.45\textwidth]{generated/logic/multiplication/multiplier-seq-cz.pdf}
\end{center}
\begin{itemize}
\item Výsledek je ve dvou registrech AC a B 
\item Pomalé, už sčítání je náročné, nyní 32 nebo 64 sčítání.
\end{itemize}
\end{frame}



\begin{frame}
\frametitle{Násobení celých čísel}

Práci sekvenční násobičky si ukážeme na příkladu:
\begin{columns}
\begin{column}{0.3\textwidth}
\texttt{\phantom{xxxx}10}\\
\texttt{\phantom{xxx}*13}\\
\vspace{-8pt}
\rule[0pt]{1.5cm}{0.1pt}\\
\texttt{\phantom{xxxx}30}\\
\texttt{\phantom{xxx}10 }\\
\vspace{-8pt}
\rule[0pt]{1.5cm}{0.1pt}\\
\texttt{\phantom{xxx}130}\\
\bigskip

\texttt{\phantom{xxxxxx}1010}\\
\texttt{\phantom{xxxxx}*1101}\\
\vspace{-8pt}
\rule[0pt]{3cm}{0.4pt}\\
\texttt{\phantom{xxxxxx}1010}\\
\texttt{\phantom{xxxxx}0000}\\
\texttt{\phantom{xxxx}1010}\\
\texttt{\phantom{xxx}1010}\\
\vspace{-8pt}
\rule[0pt]{3cm}{0.4pt}\\
\texttt{\phantom{xx}10000010}\\
\end{column}
\hfill
\begin{column}{0.4\textwidth}

\texttt{A=1010}\\
\texttt{B=1101}\\

\bigskip

\begin{center}
\includegraphics[width=0.95\textwidth]{generated/logic/multiplication/multiplier-seq-cz.pdf}
\end{center}
\end{column}
\hfill
\begin{column}{0.3\textwidth}
\begin{center}
\includegraphics[width=0.55\textwidth]{generated/logic/multiplication/multiplier-seq-example.pdf}
\end{center}
\end{column}
\end{columns}

\end{frame}


\begin{frame}
\frametitle{Rychlé násobení - Wallace tree}

Použijeme předchozí princip na co nejrychlejší součet 32 nebo 64 různých čísel:
\begin{columns}
\begin{column}{0.6\textwidth}
\includegraphics[width=1.0\textwidth]{generated/logic/multiplication/wallace-tree2.pdf}
\end{column}
\hfill
\begin{column}{0.4\textwidth}
\begin{itemize}
\item Součiny jsou jednoduché: $x_i \cdot y_j=x_i$~\texttt{and}~$y_j$
\item Nejtěžší je sečíst prostřední sloupec 64 jednobitových čísel
\item Pustíme na všechny bity, co můžeme, paralelně sčítačky a carry budeme přičítat v dalších krocích
\item V první fázi to bude 1323 sčítaček
\end{itemize}
\end{column}
\end{columns}

\end{frame}


\begin{frame}
\frametitle{Rychlé násobení - Carry Save Adder}
Jakého zrychlení tedy dosáhneme použitím sčítačky s odloženým carry?

Sčítání šesti čísel:
\begin{columns}
\begin{column}{0.3\textwidth}
\includegraphics[width=0.98\textwidth]{generated/logic/multiplication/wallace-tree6.pdf}
\end{column}
\hfill
\begin{column}{0.7\textwidth}
\begin{itemize}
\item Každý krok zredukuje velikost problému $N$ na $\lceil \frac{2}{3} \cdot N \rceil$
\item Kolik kroků potřebuje na součet 6 čísel:
\begin{itemize}
\item 1. krok redukce na součet 4 čísel - 2 zpoždění
\item 2. krok redukce na součet 3 čísel - 2 zpoždění
\item 3. krok redukce na součet 2 čísel - 2 zpoždění
\item 4. krok normální sčítačka 2 čísel - 24 zpoždění
\end{itemize}
\item Součet 6 čísel - 30 zpoždění
\item Rychlost součtu N čísel, počet kroků na redukci součtu dvou čísel $\log_{\frac{3}{2}}(N) = \frac{1}{\log(\frac{3}{2})} \cdot \log(N) = $ $ = 5.678 \cdot \log(N) = O(log(N))$
\end{itemize}
\end{column}
\end{columns}
\end{frame}



\begin{frame}[shrink=5]
\frametitle{Rychlé násobení - Wallace tree}

Když se podíváme jen na prostřední nejdelší sloupec:
\begin{center}
\includegraphics[width=0.8\textwidth]{generated/logic/multiplication/wallace-tree3.pdf}
\end{center}

\begin{itemize}
\item Vidíme, že za 8 kroků sčítaček, tedy 16 zpoždění hradel nám zbývá sečíst dva bity
\item Vpravo od tohoto sloupce je již něco sečteno a již se i propagovali carry posledních 8 sčítanců
\item Zbývá sečíst dvě 120-bitová čísla (součty a carry), což se dá stihnout asi za 30 zpoždění hradel
\item Výsledek - vynásobíme dvě čísla za cenu doby odpovídající přibližně dvěma součtům 64-bitových čísel
\item V praxi to trvá trochu déle, je použito méně sčítaček CSA.
\end{itemize}

\end{frame}



\begin{frame}
\frametitle{Dělení celých čísel}

Dělení lze zkonstruovat podobně, jako jste se učili na základní škole:
\bigskip
\begin{columns}
\begin{column}{0.4\textwidth}
\texttt{\phantom{-}120:11=10}\\
\texttt{-11}\\
\vspace{-8pt}
\rule[0pt]{1cm}{0.1pt}\\
\texttt{\phantom{xx}10}\\
\end{column}
\hfill
\begin{column}{0.4\textwidth}
\texttt{\phantom{x}01111000:1011=1010}\\
\texttt{\phantom{x}01111}\\
\texttt{\phantom{x}-1011}\\
\vspace{-8pt}
\rule[0pt]{1cm}{0.4pt}\\
\texttt{\phantom{xxx}1000}\\
\texttt{\phantom{xxx}10000}\\
\texttt{\phantom{xxx}-1011}\\
\vspace{-8pt}
\rule[0pt]{1.4cm}{0.4pt}\\
\texttt{\phantom{xxxxx}101}\\
\texttt{\phantom{xxxxx}1010}\\
\end{column}
\end{columns}
\bigskip
Oba výpočty počítají, že 120 děleno 11 je 10 a zbytek neboli \texttt{120\%11=10} binárně \texttt{1010}.

\end{frame}



\begin{frame}
\frametitle{Dělení celých čísel}

Dělička počítající \texttt{A/B}, A má 64 bitů, B 32 bitů:\\
\begin{columns}
\begin{column}{0.4\textwidth}
číslo A je uloženo ve dvou registrech AC,A\\
\includegraphics[width=1\textwidth]{generated/logic/division/divider-seq-cz.pdf}
\end{column}
\hfill
\begin{column}{0.6\textwidth}
\begin{itemize}
\item Výsledek: v registru A je podíl, v registru AC je zbytek -- modulo
\item Dělička -- existuje rychlejší algoritmus High Radix Division (je velmi složitý, přesahuje rozsah tohoto předmětu)
\begin{itemize}
\item Odhaduje několik bitů najednou a pak upřesňuje iteracemi
\item 1994 -- Pentium FDIV bug -- chyba při implementaci algoritmu Sweeney, Robertson, and Tocher (SRT) - odhaduje dva bity¨
\end{itemize}
\end{itemize}
\end{column}
\end{columns}

\end{frame}


\begin{frame}
\frametitle{Dělení celých čísel - příklad}

\begin{columns}
\begin{column}{0.35\textwidth}
\texttt{\phantom{x}01111000:1011=1010}\\
\texttt{\phantom{x}01111}\\
\texttt{\phantom{x}-1011}\\
\vspace{-8pt}
\rule[0pt]{1cm}{0.4pt}\\
\texttt{\phantom{xxx}1000}\\
\texttt{\phantom{xxx}10000}\\
\texttt{\phantom{xxx}-1011}\\
\vspace{-8pt}
\rule[0pt]{1.4cm}{0.4pt}\\
\texttt{\phantom{xxxxx}101}\\
\texttt{\phantom{xxxxx}1010}\\
\end{column}
\hfill
\begin{column}{0.4\textwidth}
číslo \texttt{A=01111000} je uloženo ve dvou registrech AC,A\\
číslo \texttt{B=1011}
\bigskip

\includegraphics[width=1\textwidth]{generated/logic/division/divider-seq-cz.pdf}
\bigskip

výsledek \texttt{A/B=\textcolor{cyan}{1010}} je uložen v registru A\\
výsledek \texttt{A\%B=1010} je uložen v registru AC

\end{column}
\hfill
\begin{column}{0.3\textwidth}
\begin{center}
\includegraphics[width=0.65\textwidth]{generated/logic/division/division-seq-example.pdf}
\end{center}
\end{column}
\end{columns}

\end{frame}


\begin{frame}
\frametitle{Rychlosti počítání}

Latence instrukcí - kolik cyklů procesoru trvá spočítat výsledek:

\begin{tabular}{l|l|l|l|l}
\textbf{Instrukce} & \textbf{Pentium M} & \textbf{IceLake} & \textbf{Zen1} & \textbf{Zen5}\\
\hline
ADD r64 & 1 (r32) & 1 & 1 & 1 \\ \hline
MUL r64 & 5 (r32)& 3 &  3 & 3\\ \hline
DIV r64 & 15-39 (r32) & 15 & 14-46 & 11-19\\ 
\end{tabular}

\url{https://www.agner.org/optimize/instruction_tables.pdf}


Délka dělení závisí na číslech, která se dělí, odhaduje se podíl několika bitů a pak se upřesňuje, což může trvat déle.

Důležitá je i propustnost, kolik oparací zvládně jedno jádro paralelně. 

Moderní CPU umožňují provést 2-6 sčítání paralelně, násobení 1-3, dělení je možno provést paralelně pouze částečně.

\end{frame}



\begin{frame}
\frametitle{Násobení a dělení čísel se znaménkem}

\begin{itemize}
\item pro násobení lze odvodit, že pokud bychom ignorovali znaménko a spočítali $M\cdot N$, pak pro $M$,$N$ obě k-bitová čísla musíme výsledek upravit takto:
\end{itemize}
$A(M\cdot N) = A(M) \cdot A(N)$\\
\phantom{$A(M\cdot N)$ }$-A(M) \cdot 2^k$ \phantom{xxxxx} pokud $N<0$\\
\phantom{$A(M\cdot N)$ }$-A(N) \cdot 2^k$ \phantom{xxxxx} pokud $M<0$
%\begin{itemize}
\begin{itemize}
\item protože reprezentace záporného čísla dvojkovým doplňkem je vlastně $A(M) = 2^k+M$, pak součin dvou záporných čísel je $(2^k+M)\cdot (2^k+N) = 2^{2\cdot k}+2^k \cdot M + 2^k \cdot N + M \cdot N$
\end{itemize}
%\end{itemize}
\begin{itemize}
\item moderní násobení a dělení počítá s absolutními hodnotami a nakonec určí znaménko výsledku podle znaménka operandů.
\begin{itemize}
\item nejvyšší bit značí znaménko
\item výpočet opačného čísla je rychlý a levný
\end{itemize}
\end{itemize}
\end{frame}


\begin{frame}
\frametitle{Kvíz - Přetečení - Overflow}

\begin{center}
\includegraphics[width=0.6\textwidth]{generated/logic/addition/overflow.pdf}
\end{center}

Obvod výše podle hodnoty nejvyšších bitů čísel v dvojkovém doplňku zjistí:
\begin{itemize}
\item[A] pouze pokud jsou oba sčítance kladné, nebo oba záporné, zda nastane přetečení.
\item[B] vždy zjistí, zda nastane přetečení
\item[C] pouze pro sčítance s rozdílnými znaménky zjistí, zda nastane přetečení
\item[D] nefunguje pro čísla v dvojkovém doplňku
\end{itemize}
\end{frame}

\begin{frame}
\frametitle{Kvíz - Přetečení - Overflow}

Může nastat přetečení v ALU při násobení a dělení čísel, v dvojkovém doplňku i bez znaménka, tak jak bylo popsáné sekvenční násobičkou a děličkou výše?

\begin{itemize}
\item[A] Může nastat přetečení při násobení i dělení
\item[B] Může nastat přetečení pouze při násobení, při dělení nemůže
\item[C] Může nastat přetečení pouze při dělení, při násobení nemůže
\item[D] Nemůže nastat ani při násobení ani při dělení.
\end{itemize}
\end{frame}

\begin{frame}
\frametitle{Kvíz - Přetečení - Overflow}

Vysvětlení:
\begin{itemize}
\item strojová instrukce násobení - násobí dva registry a výsledné číslo uloží do dvou registrů
\begin{itemize}
\item z hlediska ALU nedojde nikdy k přetečení, výsledek se do dvou registrů vejde vždy
\item program, který by uložil součin dvou int čísel do int může dojít k přetečení
\item program, který by uložil součin dvou int čísel do long long int nemůže dojít k přetečení
\end{itemize}
\item strojová instrukce dělení - dělí číslo uložené ve dvou registrech jedním registrem, výsledek uloží do jednoho registru podíl, často do druhého registru zbytek - modulo
\begin{itemize}
\item číslo složené z dvou registrů po dělení se nemusí vejít do jednoho registru, proto může dojít k přetečení
\item program, který dělí int jiným int, nemůže mít přetečení, pokud nedělí 0
\end{itemize}
\end{itemize}

\end{frame}

\section{Procesor}

\input{common/computer/tex/vonneumann-concept-cz.tex}

\section{Kódování instrukcí}

\input{common/cpu/tex/instruction-encoding-bits-cz.tex}

\input{common/cpu/tex/register-gpr-fpr-pc-riscv-cz.tex}

\input{common/cpu/tex/siglecycle-riscv-goal-cz.tex}

\input{common/cpu/tex/instruction-encoding-mips-cz.tex}

\input{common/cpu/tex/instruction-encoding-riscv-cz.tex}

\begin{frame}
\frametitle{Bonusový bod}

Jak lze v RISC-V načíst do registru konstantu o velikost 32-bitů?
\begin{itemize}
\item[A] pomocí jedné instrukce
\item[B] pomocí dvou instrukcí, první načte horních 20 bitů, druhá spodních 12-bitů
\item[C] pomocí tří instrukcí, první načte horních 16 bitů, druhá provede bitový posun doleva o 16 bitů a třetí načte spodních 16 bitů
\item[D] pomocí pěti instrukcí, první načte 12 bitů, druhá provede bitový posun doleva, třetí načte dalších 12 bitů, čtvrtá provede bitový posun doleva a pátá načte nejnižších 8 bitů
\end{itemize}
\end{frame}


\section{Konstrukce procesoru}

\input{common/logic/tex/flipflop-d-register-design-cz.tex}

\input{common/cpu/tex/cpu-cycle-pc-only-cz.tex}

\input{common/cpu/tex/cpu-building-blocks-cz.tex}

\input{common/cpu/tex/cpu-design-lw-sw-addi-alu-branch-cz.tex}

\end{document}

